% =====================================================================
%  Conferenze, Zagarolo (On The Road), giorno 1 — Talk 2
%  Entrare nell'industria a 43 anni e fare il producer per un publisher
%  Fonte: trascrizione automatica "Track 10". Talk in inglese, con slide,
%  di Tony Evans, producer di Fireshine Games, presentato al termine
%  dell'intervista del capitolo precedente. Due parti: il percorso
%  personale per entrare nell'industria; il lavoro del producer di
%  publisher, con il caso di studio di un gioco seguito per quattro anni.
%  Etichette con prefisso "c2-2:". Nessun pacchetto aggiuntivo.
% =====================================================================

\chapter[Il producer di un publisher]{Entrare nell'industria a 43 anni e fare il
producer per un publisher}
\label{cap:c2-2-producer}

Dopo il marketing (capitolo~\ref{cap:c2-1-marketing-publisher}), la stessa azienda
presenta il punto di vista della \textbf{produzione}. Il relatore è \textbf{Tony
Evans}\footnote{Così lo presenta il moderatore alla fine del capitolo precedente e
così lo ringrazia alla fine del talk (``Thanks, Tony''). In apertura la
trascrizione automatica gli attribuisce un nome diverso, chiaramente
errato.}, \textbf{producer di Fireshine Games}, gallese. Parla in inglese
(scusandosi in anticipo se andrà troppo veloce) e annuncia che resterà all'evento
tutto il fine settimana, riconoscibile da un cappellino disegnato da lui. Scherza
sul fatto che \emph{GTA VI} è appena uscito e che i presenti hanno rinunciato a
giocarci per qualche giorno.

Il talk ha due parti dichiarate:
\begin{enumerate}
  \item per chi vuole \textbf{entrare nell'industria} o si sente un po' perso: il
        suo percorso, arrivato al primo lavoro nei videogiochi \textbf{a 43 anni}
        (sezioni da~\ref{sec:c2-2-chi} a~\ref{sec:c2-2-consigli});
  \item per chi pensa di \textbf{lavorare con un publisher}: che cosa fa un
        producer di publisher, con un caso di studio (sezioni
        da~\ref{sec:c2-2-differenze} a~\ref{sec:c2-2-lezioni}).
\end{enumerate}

% ---------------------------------------------------------------------
\section{Chi è il relatore}
\label{sec:c2-2-chi}

Per ``non essere estranei'', una slide con ciò che direbbero di lui in ufficio:
\begin{itemize}
  \item è ``il gallese dell'ufficio'', come l'ospite del capitolo precedente è
        ``l'italiana'';
  \item ama i \textbf{cocktail} (non il vino, non la birra);
  \item porta \textbf{sempre un cappello} (``non perché ho perso i capelli'');
  \item tifa \textbf{Everton} (``li amo e li odio allo stesso tempo, questo è il
        calcio'');
  \item i \textbf{fumetti} lo hanno formato: si sente un supereroe ``nella misura
        in cui cerco di essere gentile, di essere buono'';
  \item ha speso centinaia di sterline nel \textbf{minigolf} (in Florida, in tre
        giorni, circa 500 euro);
  \item per lui \emph{Escape from Monkey Island} è ``il più grande videogioco mai
        creato'' (accolto da un applauso).
\end{itemize}

% ---------------------------------------------------------------------
\section{Il percorso: diciannove lavori}
\label{sec:c2-2-percorso}

Gli chiedono spesso come sia entrato nell'industria, perché ha avuto il primo
lavoro \textbf{a 43 anni}, ``dieci giorni prima del 44º compleanno''. Spera che il
racconto arrivi a chi, tra i più giovani, sta per arrendersi.

Preparando le slide ha contato \textbf{19 lavori}: commercio, sanità, teatro,
pubblica amministrazione, enti di beneficenza, istruzione. Ha lavorato anche a
\textbf{Legoland}, nel turno di notte: distruggeva le costruzioni fatte dai bambini
durante il giorno (``ho ancora le cicatrici sui pollici''). Il punto: \textbf{si
imparano cose lungo la strada}, senza rendersene conto, e le si porta con sé.

\subsection{La scuola di cinema}

I videogiochi non erano nei suoi piani: voleva fare il \textbf{regista} e ha fatto
la scuola di cinema. Ha imparato che \textbf{non voleva fare film} (``lo odiavo
davvero''), ma anche che \textbf{pianificare} un film (logistica, cast, troupe,
organizzazione) gli riusciva bene e gli piaceva. Dopo la scuola, entrare nel cinema
era difficile quanto nei videogiochi; ha dovuto pagare le bollette, è stato
promosso, si è adagiato in lavori amministrativi, fino a un licenziamento e poi a
Legoland.

\subsection{La National Film and Television School}

Sua moglie seguiva un corso alla \textbf{NFTS} e gli segnalò un posto: divenne
\textbf{coordinatore del corso di videogiochi} della scuola. Era il 2013 circa, e
paradossalmente fu Legoland a fargli ottenere il posto: lavorando dalle 23 alle 7
ascoltava \textbf{podcast sui videogiochi}, nel periodo in cui esplodevano gli
indie, e al colloquio poté parlare di \emph{Journey}, \emph{Spelunky}, \emph{Fez}:
esattamente ciò che il responsabile del corso voleva.

A cambiarlo non furono la scuola o il lavoro, ma \textbf{le persone}: studenti e
docenti. Aveva voce in capitolo su chi invitare a parlare agli studenti
(animazione, programmazione, produzione). Un giorno, chiacchierando con un
producer, gli descrisse le sue giornate e quello rispose: ``ma questo è essere un
producer''. Fu l'\textbf{epifania}: ``questo è un lavoro che posso fare''. In sei
anni e mezzo alla NFTS ha aiutato circa \textbf{50 persone} a entrare
nell'industria, e molti studenti sono diventati amici. ``Non posso sottolineare
abbastanza'' quanto quelle connessioni lo abbiano portato dov'è.

\subsection{Lo studio mai nato}

Nel 2019 decise che ``forse era il suo turno''. Tornò in Galles per fondare uno
studio con un'idea che considera ancora ``geniale'': nei corsi gli studenti
\textbf{raramente finiscono i loro giochi}, che però sono spesso buoni; lo studio
li avrebbe presi e completati. Aveva persone a bordo, un business plan, contatti
con il governo gallese. Poi arrivò il \textbf{Covid}: ``il denaro è sparito
schioccando le dita''.

Per pagare le bollette prese un lavoro all'università locale: ``il peggior lavoro
della mia vita'', ``ogni giorno una roulette russa a seconda dell'umore della
responsabile''. La lezione: \textbf{una sola persona può cambiare l'umore di un
intero team}; in quel caso in negativo. ``Non siate quella persona.'' Lui cerca di
essere quello rilassato, ``il simpatico della stanza''. Un'altra lezione: cavare
qualcosa di buono da una situazione brutta. Organizzò per il suo team una giornata
nel laboratorio del corso di \textbf{scienze forensi}, con tute e una finta scena
del crimine. Poi, con il Covid, fu di nuovo licenziato.

\subsection{Fireshine, ``dal nulla''}

Fireshine arrivò ``di punto in bianco'': \textbf{furono loro a contattarlo}. Tutto
risaliva a una conoscenza: nel 2016 aiutava alcuni ex studenti con il loro studio
(fogli di calcolo, budget, contatti con i publisher). Lo studio partecipava a un
programma in Irlanda del Nord che assegnava un \textbf{mentore di produzione},
Neil. I due si trovarono bene, si rividero negli anni agli eventi, e Neil,
probabilmente seguendo su LinkedIn le sue difficoltà a trovare lavoro, \textbf{lo
propose a Fireshine senza dirglielo}.

Il colloquio, per un posto da \textbf{scout}, fu ``il peggiore della mia vita'':
dopo venti minuti lo interruppe chiedendo ``è inutile, vero?''. Lo era, perché non
lo volevano come scout: lo stavano valutando per un altro ruolo, e divenne il
\textbf{primo release manager} di Fireshine. Ciò di cui è grato: \textbf{non hanno
guardato i suoi crediti} (non ne aveva) \textbf{ma le sue competenze}. Dall'epifania
al primo lavoro aveva inviato ``facilmente \textbf{500 candidature}'' (forse il
doppio) e ottenuto forse sei colloqui: senza titoli pubblicati è durissima.

Dopo circa un anno come release manager seppe che cercavano un \textbf{associate
producer}; lo disse al capo e il giorno dopo lo era. ``A volte basta dire ciò che
hai in testa.'' Ma ciò che aiutò furono i dodici mesi precedenti: lo conoscevano,
conoscevano la sua etica del lavoro e l'impegno con cui aveva costruito un ruolo
nuovo, imparando da zero i \emph{back end} di Sony, Microsoft e Nintendo (``un
campo minato''), facendo molti errori (``che non vi dirò mai, perché loro ancora
non li conoscono'').

% ---------------------------------------------------------------------
\section{Consigli per entrare}
\label{sec:c2-2-consigli}

\begin{description}
  \item[Non essere uno stronzo.] La regola gallese che ripete spesso
        (\emph{don't be a twat}). Alla NFTS, durante i colloqui di ammissione,
        faceva ``il test'': la persona che si presenta ai futuri docenti è diversa
        da quella che si incontra accompagnandola in aula. Lo stesso vale con gli
        studi: per lui e per l'azienda è più importante essere \textbf{persone con
        cui è facile lavorare}. Un successo è bello, ma se i due anni di lavoro sono
        un incubo nessuno è felice. Non farsi mettere i piedi in testa, sapere che
        cosa si vuole, ma con serenità. Fireshine ha \textbf{rifiutato giochi poi
        diventati grandi successi} perché non era sicura del team.
  \item[Fai domande, sii curioso.] Ma c'è una linea sottile tra domande fatte per
        il gusto di farle e domande \textbf{genuine}: la gente se ne accorge.
  \item[Mostra, non raccontare.] Dimostrare interesse partecipando: \emph{game
        jam}, ma anche offrire aiuto a uno studio che ha un problema di
        produzione, come fece lui. E ancora: i \textbf{playtest}, perché tutti
        hanno bisogno di qualcuno che provi i loro giochi.
  \item[Fai networking.] Lo odia, è ``un'enorme debolezza''; agli eventi lo aiuta
        rappresentare l'azienda, ``una coperta di sicurezza''. Quella mattina ci
        ha messo quindici minuti per presentarsi a un gruppo di relatori. Ma agli
        eventi \textbf{è previsto} che la gente si avvicini e parli: non abbiate
        paura.
  \item[Usa LinkedIn.] Anche LinkedIn può dare coraggio, come un avatar: scrivere
        qualcosa e andarsene è più facile che parlare di persona. Lui lo usa come
        una rubrica e ha oltre 2.000 contatti.
  \item[Le raccomandazioni contano.] Dei suoi 19 lavori, \textbf{nove} sono
        arrivati grazie a una raccomandazione. Le aziende rischiano più volentieri
        su una persona raccomandata che su uno sconosciuto. E quando si è dentro,
        bisogna \textbf{restituire}: lui stesso ha fatto entrare qualcuno in
        Fireshine.
  \item[Non succederà in una notte.] Candidarsi diventa un lavoro a tempo pieno;
        la pazienza è la chiave. Basta \textbf{una persona} che guardi le tue
        competenze e ti veda.
  \item[Sii tattico.] Non voleva fare il release manager né lo scout: voleva
        \textbf{entrare nell'industria}. A volte entrare è l'obiettivo, anche se si
        punta più in alto: lui ha impiegato quattro anni a diventare producer, e in
        Fireshine molti hanno cambiato ruolo (da product manager a community
        manager, business development, influencer manager).
  \item[Non è una gara.] Non preoccuparsi degli altri: sì, a volte chi ti passa
        davanti lo fa ingiustamente (``magari il padre è nell'industria''), ma
        pensarci fa impazzire.
\end{description}

Il messaggio chiave: ``non sono nessuno'', in senso positivo. Non aveva un dono
nascosto: \textbf{era un producer da sempre, senza saperlo}. È ``una persona
normale'' che ha fatto molti lavori e raccolto molte competenze che all'industria
servivano: \textbf{soft skill}, saper parlare e comunicare con le persone, saper
usare un foglio di calcolo e, soprattutto, \textbf{buon senso}. ``Se ce l'ho fatta
io, potete farcela voi.''

Oggi è producer da circa quattro anni (dopo un periodo da associate più lungo di
quanto avrebbe voluto) e ama il suo lavoro, che sfrutta le sue doti naturali e
soddisfa anche il prurito creativo del cinema. In cinque anni e mezzo in Fireshine
``si contano sulle dita di una mano'' le persone che se ne sono andate: l'azienda
cresce ma la cultura resta la stessa, e l'azienda ``si è guadagnata la sua
lealtà'' perché lo ha ascoltato e si è fidata di lui. Tra i titoli principali del
publisher cita \emph{Core Keeper}, \emph{Far Far West} e un gioco
\textbf{candidato come miglior gioco indipendente ai The Game Awards}, la prima
candidatura di Fireshine\footnote{Il titolo è trascritto ``DentroTac'', ``Denture
Attack'', ``Bench Attack'': non è stato possibile identificarlo.}.

% ---------------------------------------------------------------------
\section{Producer di studio e producer di publisher}
\label{sec:c2-2-differenze}

Che differenza c'è? Tony non ha mai lavorato in uno studio, ma ha chiesto in giro e
avverte che le cose cambiano da studio a studio.
\begin{center}
\begin{tabular}{@{}p{0.44\textwidth}p{0.44\textwidth}@{}}
\toprule
\textbf{Producer di studio} & \textbf{Producer di publisher} \\
\midrule
Workflow, scadenze, tracciamento dei bug, deliverable &
Approccio più ampio, senza le minuzie quotidiane \\
Integrato nel lavoro quotidiano del team (``la parte che invidio'') &
Una chiamata settimanale fissa con ogni team (di più se le cose vanno male),
sempre disponibile su Discord \\
Scrum e processi, a seconda della composizione dello studio &
Più progetti insieme (lui ne ha quattro): priorità e giocoleria \\
 & \textbf{Risponde al business}: allineare gli obiettivi del dev, del publisher
e della proprietà del publisher \\
\bottomrule
\end{tabular}
\end{center}

\subsection{Che cosa facciamo, secondo i colleghi e i dev}

Spiegarlo è difficile: alla famiglia dice ``project manager''. Per gioco ha chiesto
ai colleghi una definizione:
\begin{itemize}
  \item ``\textbf{Proteggo le persone intelligenti dalle distrazioni}'' (Jamie);
  \item ``Make game made good'' (Cormac, con ironia);
  \item ``Aiutiamo i dev a portare il gioco al traguardo, \textbf{tappando i buchi}
        se il team non ha un producer dedicato'' (Jeremy).
\end{itemize}
Ha chiesto anche a due sviluppatori che segue: entrambi evidenziano
\textbf{organizzazione}, \textbf{logistica} ed essere \textbf{il punto di
contatto} con Fireshine, un canale verso i vari reparti del publisher.

\subsection{La sua definizione}

Ispirata a Batman: \textbf{``sono ciò di cui il team ha bisogno''}. Ogni team è
diverso, e lui diventa un producer diverso a seconda del gioco e delle persone. I
suoi quattro progetti:
\begin{description}
  \item[\emph{Shadows of Doubt}.] Un \emph{detective sim} generato proceduralmente,
        in voxel art, uscito da un paio d'anni; sviluppato \textbf{da una sola
        persona} (Cole). Qui Tony si concentra su scadenze, QA e porting, perché le
        patch funzionino su console come su PC.
  \item[\emph{Food Truck Empire}.] Un gestionale, il team più grande (12 persone,
        a Varsavia), già esperto nel pubblicare giochi: anche qui scadenze e
        test, con un buon ritmo di playtest.
  \item[\emph{Husk Protocol}.] Un \emph{deck builder} con mecha, team di tre
        persone con trent'anni di esperienza (Ubisoft, 11 bit studios), ma al primo
        gioco come studio: riceve ``un po' di tutto''. È il più vicino all'uscita
        (primo trimestre dell'anno successivo), quindi le priorità si stanno
        spostando su di lui.
  \item[\emph{Rekka}.] Il caso di studio del talk.
\end{description}

% ---------------------------------------------------------------------
\section{Il caso di studio: Rekka}
\label{sec:c2-2-rekka}

\emph{Rekka}\footnote{Trascritto anche ``Wrecker'', ``Brecker'', ``Recca'',
``Racket''. Lo studio è trascritto ``Ampersform Entertainment'' ed
``Emberstall''; la grafia esatta non è verificabile.} è sviluppato da uno studio di
\textbf{Berlino} di cinque persone (``tre sono altissime''), \textbf{al primo
gioco}:
\begin{itemize}
  \item firmato nel \textbf{2022};
  \item in \textbf{accesso anticipato} dal settembre 2024;
  \item uscito in \textbf{1.0} poche settimane prima del talk;
  \item poco meno di \textbf{300.000 copie} vendute, il sesto titolo più venduto di
        Fireshine;
  \item \textbf{oltre un milione di wishlist}. ``Come trasformarle in acquisti? Se
        lo sapessimo, avremmo risolto i videogiochi.''
\end{itemize}
Tony lo segue da quattro anni, da prima della firma. Non è stato facile: malattie,
ritardi, e l'accesso anticipato non era nemmeno il piano originale. ``Molte
improvvisazioni.''

\subsection{Le attività del producer di publisher}

Una slide elenca i compiti in tre colonne: \textbf{interni} (per il business di
Fireshine), \textbf{sul gioco} (qualità del prodotto), \textbf{con i partner}
(organizzazione e rapporti esterni). Quando uno studio firma vuole delegare molto,
e Fireshine è felice di farlo; nelle prime riunioni chiede che cosa lo studio vuole
fare da sé (per esempio gestire la community), ma su certe cose ``non diamo nemmeno
la scelta'': le prende il publisher perché richiedono molto tempo.

\subsection{QA}

Su \emph{Rekka} soprattutto \textbf{QA funzionale}, a fasi alterne per quasi
quattro anni, poi circa \textbf{tre mesi continuativi} prima della 1.0. Il QA era
affidato a Lionbridge, a Varsavia: Tony è andato a conoscere i tester, perché era
importante che \textbf{si sentissero parte del team} (``se si sentono parte del
team faranno di più''). Al lead tester, che ha seguito il gioco per tre anni su
quattro, ha regalato una maglietta del gioco: ``piccole cose''.

\subsection{Localizzazione}

Molte lingue (inglese, tedesco, giapponese, cinese, ucraino, polacco, russo\dots;
``niente italiano, scusate''). La traduzione non è \textbf{parola per parola}:
altrimenti basterebbe Google Translate. È questione di \textbf{sfumature}. Per un
gioco che è \textbf{una fiaba}, una storia da tramandare di generazione in
generazione, il \textbf{tono} era fondamentale: il testo inglese era bellissimo e
``lo dovevamo'' alle altre lingue. Tony ha dovuto \textbf{cambiare fornitore}
perché alcuni testi non lo soddisfacevano.

\subsection{Quanto costa}

``Visto che non ci sono telecamere'', mostra i costi: il QA funzionale arriva a
circa \textbf{115.000}, il totale di QA e localizzazione a circa \textbf{200.000
euro}. Può sembrare molto, ma va contestualizzato: quattro anni, molte demo, molti
aggiornamenti in accesso anticipato, ognuno da testare. Il bisogno di test varia:
\emph{Husk Protocol} probabilmente ne richiederà meno, anche per come il team ha
impostato codice e ambiente in Unity fin dall'inizio. Il problema di \emph{Rekka}
era che doveva essere \textbf{un gioco piccolo}, impostato come tale, e poi è
cresciuto: Tony lo immagina come ``un triangolo rovesciato'', con tutto che si
riversa su una base piccola. Visto il lavoro degli ultimi nove mesi per renderlo
degno dell'uscita, avrebbe speso anche di più, ``e sarei stato comunque
tranquillo'': la 1.0 è uscita con recensioni molto positive.

Un altro costo nascosto: con uno \emph{writer} a bordo che produce molto testo, è
difficile essere disciplinati; capita di \textbf{tradurre testo che poi non entra
nel gioco}, e quei soldi non tornano.

\subsection{Porting e user testing}

\begin{itemize}
  \item \textbf{Porting}: zero su \emph{Rekka}, per ora. Se ne riparlerà in base
        alle vendite; ``Tony vorrebbe la Switch'', ma serve un business case:
        ``non buttiamo soldi su qualcuno per vedere che cosa succede''. Su
        \emph{Shadows of Doubt} invece le patch per console sono parte del lavoro,
        e altri producer di Fireshine sono ``molto dentro'' il porting (un titolo
        del catalogo è arrivato anche su mobile).
  \item \textbf{User testing}: su \emph{Rekka} solo due sessioni, prima di
        Gamescom, per verificare che il \textbf{tutorial} fosse comprensibile.
        Oggi lo strumento è molto evoluto: per \emph{Husk Protocol} si ottiene un
        \textbf{video di gioco con commento}, una \textbf{trascrizione} con
        etichette per saltare ai punti rilevanti e un riassunto marcato
        chiaramente come \textbf{generato dall'IA}. Tony insiste che i dev guardino
        i video e si facciano un'idea propria: ``l'IA a volte esagera''; è una buona
        guida, non una risposta definitiva.
\end{itemize}

\subsection{Feedback e playtest}

\begin{itemize}
  \item Un \textbf{playtest aperto su Steam} con \textbf{1.750 giocatori}: forse
        pochi, ma con un team piccolo è difficile gestire troppo feedback.
  \item Il \textbf{feedback interno} si adatta al team: i dev di \emph{Rekka}
        volevano i suoi appunti su Discord ``come flusso di coscienza'' mentre
        giocava; il team di \emph{Food Truck Empire} vuole video con commento.
        ``Sono ciò di cui il team ha bisogno.''
  \item Dal feedback e dai messaggi sui social era chiaro che c'era \textbf{una
        community investita nel gioco}: di qui la decisione, ``non facile né
        veloce'', di passare all'\textbf{accesso anticipato}. Il pubblico ha
        contribuito molto a dare forma al gioco.
  \item Due strumenti di tracciamento: una \textbf{bacheca} con bug e feedback dei
        giocatori (``la più caotica che abbia mai visto, ma funziona per i dev,
        quindi faccio un passo indietro''); e uno strumento che \textbf{estrae le
        recensioni di Steam} in un foglio, indicando se positive o negative e se
        il gioco è stato rimborsato.
  \item Nei giorni prima della 1.0, Tony e il team community hanno
        \textbf{risposto a ogni singola recensione negativa}: se il problema era
        risolto, lo dicevano e invitavano a riprovare. ``Mi ci è voluto il
        weekend'', ma diverse recensioni sono state ribaltate.
  \item \textbf{Beta della community}: sul Discord (circa 13--15.000 membri)
        chiunque poteva giocare build non definitive in cambio di feedback, per
        tutto l'accesso anticipato. Per la 1.0, temendo gli spoiler sul finale,
        hanno fatto una \textbf{beta chiusa} con circa 40 persone per verificare
        l'impatto emotivo della storia. ``Il feedback è la linfa vitale.''
\end{itemize}

\subsection{Roadmap, risorse e la sceneggiatrice}

La \textbf{roadmap} è ``il pane quotidiano'' di ogni chiamata con i dev: contiene
i compiti del producer e i \textbf{deliverable} del team, ciascuno con una soglia
di qualità da superare per l'approvazione.

Non tutti gli studi hanno contatti, soprattutto se appena usciti da una scuola; il
publisher li fornisce. Per \emph{Rekka}:
\begin{itemize}
  \item uno studio di \textbf{co-sviluppo} per qualche mese di supporto alla
        programmazione e un altro per l'animazione;
  \item una sceneggiatrice: il team voleva ``qualcuno come Rhianna
        Pratchett''\footnote{Trascritto ``Rihanna Pratchett''.}, e Tony ha
        portato \textbf{Rhianna Pratchett}. Il nome fa notizia, e soprattutto ha
        mostrato al team che il publisher ``non scherza'': ``se dite che volete una
        cosa, cerchiamo di ottenerla''. L'unico rammarico: era così brava che si
        voleva usare tutto, e non si poteva.
\end{itemize}

\subsection{I compiti interni}

\begin{description}
  \item[Budget.] Su \emph{Rekka} l'uscita è arrivata circa \textbf{due anni dopo}
        il previsto, per una scelta di business legata all'accesso anticipato e
        alla community; eppure il progetto si è chiuso \textbf{sotto budget}.
  \item[Rischio.] La gestione del rischio è quotidiana, e di solito ruota intorno a
        \textbf{qualità, tempo e denaro}: ``questa build ha dei bug; per averne
        meno servono tanto tempo e tanti soldi; ne vale la pena?''. Non è garanzia
        di successo, e le recensioni negative arriveranno comunque: ``abituatevi ai
        bug''.
  \item[Milestone.] Ogni ritardo cambia il contratto: Tony ha \textbf{riscritto
        le milestone di Rekka sette volte}. Le definisce insieme ai dev
        \textbf{prima della firma}, perché siano raggiungibili: ``non cerco di
        fregare nessuno, voglio pagare i dev''. Dal lavoro nel volontariato ha
        imparato che \textbf{i mezzi di sussistenza delle persone dipendono dai
        pagamenti}.
  \item[Stakeholder.] Tutti in azienda toccano il gioco: produzione, marketing,
        community, vendite, business development, perfino l'amministrazione che
        paga le fatture. Tony ha spinto perché \textbf{tutti vengano accreditati},
        anche per una conversazione che ha fatto nascere un'idea. La comunicazione
        interna passa da un canale Teams; il collega del business development, ben
        informato, ha trovato per esempio un co-publisher.
  \item[Eventi.] Gli eventi servono a dare un volto ai nomi. A una fiera ha
        conosciuto qualcuno che poi ha fatto feedback sui loro giochi. La prima
        volta di \emph{Rekka} a Gamescom c'era \textbf{una coda di tre ore}, tanto
        da dover limitare le partite a venti minuti (``mi odiavano'').
\end{description}

% ---------------------------------------------------------------------
\section{Processo, prodotto, persone}
\label{sec:c2-2-persone}

I producer, dice Tony anticipando un talk del giorno dopo
(capitolo~\ref{cap:c3-3-people-over-processes}), tendono a ricadere in tre
categorie: \textbf{processo}, \textbf{prodotto}, \textbf{persone}. Lui è decisamente
uno da \textbf{persone}, e pensa che gli abbiano affidato \emph{Rekka} proprio per
questo: anni passati con gli studenti, un carattere rilassato e avvicinabile,
``probabilmente non ciò che la gente si aspetta da un \emph{publisher
malvagio}''.

Per quattro anni ha parlato con il team \textbf{ogni martedì alle 11:30}; ora che
il gioco è uscito i contatti si diradano, ``ed è strano''. Cerca di
\textbf{integrare} i dev nel team Fireshine: nei nuovi uffici le sale riunioni
portano i nomi dei giochi, e c'è una ``sala Rekka''. Al momento del lancio della
1.0 ha voluto che \textbf{tutti premessero insieme il pulsante} di pubblicazione:
per lui è cliccare qualche tasto, per loro è una pietra miliare personale e di
studio. E mostra il tatuaggio di \emph{Rekka} sul braccio: dentro un publisher ci
sono molte persone a cui il gioco \textbf{importa davvero}.

La sintesi: ha guadagnato la loro \textbf{fiducia} facendo ciò che diceva; la loro
opinione restando \textbf{autentico}; il loro ascolto essendo \textbf{schietto}
(``non li ho mai presi in giro''); e ha mantenuto l'impegno del primo giorno:
restare con loro fino alla fine.

% ---------------------------------------------------------------------
\section{Lezioni da producer di publisher}
\label{sec:c2-2-lezioni}

\begin{enumerate}
  \item \textbf{Le visite in studio rendono umani.} Andare spesso a Berlino ha
        messo il team a suo agio nel parlargli di tutto, anche delle cose personali
        e delle conversazioni difficili.
  \item \textbf{A volte non fare nulla è fare qualcosa.} È forte la tentazione di
        essere il cavaliere che salva la situazione; ma \textbf{i dev conoscono il
        gioco}. Spesso basta fare un passo indietro e osservare: di solito risolvono
        da soli. Se non succede, si interviene.
  \item \textbf{Togliti di mezzo, apri la strada.} Rendere il più facile possibile
        ai dev ``fare i dev'': togliere compiti, offrire aiuto (ora sta riscrivendo
        alcuni tutorial).
  \item \textbf{Produci il team, non il gioco.} Gran parte del lavoro è far
        funzionare il team al meglio, anche con un \textbf{coaching} fatto ``per
        osmosi''. Anche chi ha decenni di esperienza a volte va aiutato a cambiare
        prospettiva, soprattutto sul \textbf{feedback dei giocatori}: è facile dire
        ``è una sciocchezza'', ma dietro ogni feedback c'è una ragione, che vale la
        pena capire anche se poi lo si scarta.
  \item \textbf{Fidati dell'istinto, ma anche dei dati.} Il suo istinto ha avuto
        quasi sempre ragione, anche quando non l'ha seguito. Il suo grande
        rimpianto è aver lasciato che \textbf{l'anzianità altrui} prevalesse sulle
        sue opinioni, da nuovo arrivato in stanze piene di direttori: ``avevo
        ragione, e vorrei aver parlato di più''. Non basta fare le domande giuste:
        bisogna \textbf{interrogare le risposte}.
  \item \textbf{È buon senso.} Un veterano con cui lavorava (Graham), a ogni
        domanda su come sapesse che cosa serviva, rispondeva ``è buon senso'', e
        lo faceva impazzire. Ora Tony ha capito: quando hai le informazioni, conosci
        le persone e le conseguenze, decidere è facile.
  \item \textbf{Non sai ciò che non sai.} Non bisogna avere paura di dire ``non lo
        so'': \textbf{bluffare con un team fa perdere la fiducia}. E non si può
        pianificare tutto: la notte prima, alle 2, stava copiando le note sul tablet
        perché la presentazione non funzionava. Improvvisare e pensare in fretta:
        ``l'ho imparato alla scuola di cinema''.
\end{enumerate}

% ---------------------------------------------------------------------
\section*{Punti chiave del capitolo}
\addcontentsline{toc}{section}{Punti chiave del capitolo}

\begin{itemize}
  \item Si può entrare nell'industria tardi: 19 lavori, 500 candidature, il primo
        posto a 43 anni. Le \textbf{competenze trasferibili} (organizzazione,
        soft skill, buon senso) valgono quanto i crediti.
  \item Le \textbf{connessioni} sono decisive: nove lavori su diciannove da
        raccomandazioni; il posto in Fireshine da un mentore conosciuto anni prima.
  \item Consigli: non essere uno stronzo (i publisher rifiutano team difficili),
        domande genuine, mostrare invece di raccontare, networking anche via
        LinkedIn, pazienza, essere tattici (entrare prima, crescere poi), non
        confrontarsi con gli altri.
  \item Il producer di publisher segue più progetti, con chiamate settimanali, e
        \textbf{risponde al business}; è il punto di contatto con il publisher e
        ``ciò di cui il team ha bisogno''.
  \item Compiti: QA, localizzazione (tono e sfumature, non parola per parola),
        porting, user testing (l'IA come guida, non come verità), playtest,
        community beta, risposta alle recensioni negative, roadmap, risorse
        esterne (co-sviluppo, sceneggiatori), budget, rischio
        (qualità--tempo--denaro), milestone, stakeholder, eventi.
  \item Le milestone si costruiscono con i dev prima della firma: dai pagamenti
        dipendono le persone.
  \item Lezioni: visite in studio, a volte non intervenire, togliere ostacoli,
        produrre il team, ascoltare il perché del feedback, fidarsi dell'istinto e
        dei dati, non bluffare.
\end{itemize}
